\section{Sobolev范数与完备性}
\label{sec:sobolev-norms-completeness}

上一节的例子要求把函数与它的弱导数放在一起度量。本节固定三项数据，即区域$\Omega$、导数的最高阶数$k$和积分指数$p$，用$L^p$范数分别度量函数及各阶弱导数，再合成一个范数。然后证明这样得到的空间完备，并在$p=2$时写出它的内积和坐标形式。

\subsection{空间的定义}
\label{sec:sobolev-space-definition}

先约定记号。在欧氏开集$\Omega\subset\mathbb R^d$上，多重指标$\alpha=(\alpha_1,\ldots,\alpha_d)$由非负整数组成，$|\alpha|=\alpha_1+\cdots+\alpha_d$是导数的总阶数。$\partial^\alpha f$表示第~\ref{chap:smoothing-weak-derivatives}章意义下的弱导数，约定$\partial^0f=f$。

要询问$L^p$函数是否有弱导数，先要知道它局部可积。对紧集$K\subset\Omega$和$p<\infty$，Hölder不等式给出
\[
\int_K|f|\,\mathrm dx\leq|K|^{1-1/p}\lVert f\rVert_p,
\]
$p=\infty$时则有$\int_K|f|\leq|K|\,\lVert f\rVert_\infty$。所以$L^p$函数都局部可积，检验弱导数总是有意义的。

函数与每个导数各占一项，用下面的范数合在一起。

\begin{definition}[Sobolev空间]
\label{def:sobolev-space}
设$\Omega\subset\mathbb R^d$为开集，$k\geq0$为整数，$1\leq p\leq\infty$。所有满足下述条件的$f\in L^p(\Omega)$组成的集合记作$W^{k,p}(\Omega)$：对每个$|\alpha|\leq k$，弱导数$\partial^\alpha f$存在且属于$L^p(\Omega)$。称它为\term{Sobolev空间}（Sobolev space），并赋以\term{Sobolev范数}（Sobolev norm）
\begin{equation}
\lVert f\rVert_{W^{k,p}(\Omega)}=
\begin{cases}
\Bigl(\sum\limits_{|\alpha|\leq k}\lVert\partial^\alpha f\rVert_{L^p(\Omega)}^p\Bigr)^{1/p}, & 1\leq p<\infty,\\
\max\limits_{|\alpha|\leq k}\lVert\partial^\alpha f\rVert_{L^\infty(\Omega)}, & p=\infty .
\end{cases}
\label{eq:sobolev-norm}
\end{equation}
\end{definition}

三项数据各有分工。阶数$k$规定控制到第几阶导数，指数$p$规定函数和每个导数用哪种积分方式度量大小，区域$\Omega$规定在哪里检验测试恒等式。$k=0$时没有导数要求，$W^{0,p}(\Omega)=L^p(\Omega)$。区域不致混淆时，把$\lVert\cdot\rVert_{L^p(\Omega)}$简写为$\lVert\cdot\rVert_p$。

式~\eqref{eq:sobolev-norm}确实是范数。齐次性直接成立。三角不等式分两步：先对每个分量用$L^p$的Minkowski不等式，得到
\[
\lVert\partial^\alpha(f+g)\rVert_p\leq\lVert\partial^\alpha f\rVert_p+\lVert\partial^\alpha g\rVert_p,
\]
再把各分量的范数排成一个有限维向量，用$\ell^p$的Minkowski不等式合并。若范数为零，则零阶项$\lVert f\rVert_p=0$，$f$是零等价类。

把每个导数都计入范数，作用是让求导成为连续线性映射。对$|\alpha|\leq k$，把$\partial^\alpha$看成从$W^{k,p}(\Omega)$到$L^p(\Omega)$的映射，直接得到
\begin{equation}
\lVert\partial^\alpha f-\partial^\alpha g\rVert_p\leq\lVert f-g\rVert_{W^{k,p}} .
\label{eq:sobolev-derivative-continuous}
\end{equation}
这并不说明求导在$L^p$上连续：第~\ref{sec:sobolev-oscillation}小节的$f_n$在$L^2$中大小不变，导数却无界。连续性来自定义域的选择，导数的大小已经预先计入了输入的范数。

这里可以再用一次向量的类比。把$f$对应到有限个分量组成的“长向量”$(\partial^\alpha f)_{|\alpha|\leq k}$，它是乘积空间$L^p\times\cdots\times L^p$中的点，Sobolev范数就是这个长向量的范数。不同的是，这些分量并不独立：它们必须是同一个函数的各阶导数，受测试恒等式约束。Sobolev空间是乘积空间中满足这组约束的子集。

空间的元素是几乎处处等价类，弱导数也只确定到几乎处处。选不同的代表不改变元素，也不改变它的弱导数，但“某一点的函数值”一般没有意义，后面讨论端点值和边界值时必须专门处理。判断$f\in W^{k,p}(\Omega)$时，须对测试函数核对积分恒等式，不能用几乎处处的经典导数代替。

\begin{example}[折点函数的Sobolev范数]
\label{ex:sobolev-kink-norm}
在$(0,1)$上，$f(x)=|x-1/2|$的弱导数是$\operatorname{sgn}(x-1/2)$，二者都有界，所以$f\in W^{1,p}(0,1)$对每个$1\leq p\leq\infty$成立。

$p=2$时，由第~\ref{sec:sobolev-kinks-jumps}小节的计算，
\[
\lVert f\rVert_{W^{1,2}}^2=\frac1{12}+1=\frac{13}{12}.
\]
$p=\infty$时，两项的上确界分别是$1/2$和$1$，所以
\[
\lVert f\rVert_{W^{1,\infty}}=\max\Bigl(\frac12,\,1\Bigr)=1.
\]

$f$不属于$W^{2,p}(0,1)$。若二阶弱导数$h$存在，则$h$是$\operatorname{sgn}(x-1/2)$的弱导数。而
\[
\operatorname{sgn}(x-1/2)=2\,\mathbf 1_{[1/2,1)}(x)-1\quad\text{（几乎处处）},
\]
所以$h/2$是阶跃函数$H$的弱导数，这与上一节的结论矛盾。
\end{example}

\subsection{完备性}
\label{sec:sobolev-completeness}

$L^p(\Omega)$的完备性已在第~\ref{chap:function-normed-spaces}章证明，所以长向量$(\partial^\alpha f)_\alpha$的每个分量都有极限。需要说明的是，取极限以后各分量之间的约束仍然成立，也就是极限点仍在这个子集中。在乘积空间的图景里，这就是说该子集是闭的。

\begin{theorem}[Sobolev空间的完备性]
\label{thm:sobolev-complete}
设$\Omega\subset\mathbb R^d$为开集，$k\geq0$为整数，$1\leq p\leq\infty$。则$W^{k,p}(\Omega)$是Banach空间。
\end{theorem}

\begin{proof}
设$(f_n)$是$W^{k,p}(\Omega)$中的Cauchy列。

\emph{各分量的极限。}由式~\eqref{eq:sobolev-derivative-continuous}，对每个$|\alpha|\leq k$，$(\partial^\alpha f_n)_n$都是$L^p(\Omega)$中的Cauchy列。由$L^p$的完备性，存在$g_\alpha\in L^p(\Omega)$，使
\[
\partial^\alpha f_n\to g_\alpha\quad\text{在 }L^p(\Omega)\text{ 中}.
\]
记$f=g_0$。此时还不知道$g_\alpha$与$f$有什么关系，它们只是分别得到的极限。

\emph{把测试恒等式传到极限。}固定$\varphi\in C_c^\infty(\Omega)$，对每个$n$有
\[
\int_\Omega\partial^\alpha f_n\,\varphi\,\mathrm dx=(-1)^{|\alpha|}\int_\Omega f_n\,\partial^\alpha\varphi\,\mathrm dx .
\]
设$p'$为$p$的共轭指数。$\varphi$与$\partial^\alpha\varphi$有界且支撑紧，所以都属于$L^{p'}(\Omega)$。由Hölder不等式，
\[
\begin{aligned}
\Bigl|\int_\Omega(\partial^\alpha f_n-g_\alpha)\varphi\,\mathrm dx\Bigr|
&\leq\lVert\partial^\alpha f_n-g_\alpha\rVert_p\lVert\varphi\rVert_{p'},\\
\Bigl|\int_\Omega(f_n-f)\partial^\alpha\varphi\,\mathrm dx\Bigr|
&\leq\lVert f_n-f\rVert_p\lVert\partial^\alpha\varphi\rVert_{p'}.
\end{aligned}
\]
两式右端都趋于零。令$n\to\infty$，得到
\[
\int_\Omega g_\alpha\varphi\,\mathrm dx=(-1)^{|\alpha|}\int_\Omega f\,\partial^\alpha\varphi\,\mathrm dx
\]
对所有测试函数成立，即$g_\alpha=\partial^\alpha f$。

\emph{合并分量。}所有$\partial^\alpha f=g_\alpha$都属于$L^p(\Omega)$，所以$f\in W^{k,p}(\Omega)$。$p<\infty$时，
\[
\lVert f_n-f\rVert_{W^{k,p}}^p=\sum_{|\alpha|\leq k}\lVert\partial^\alpha f_n-g_\alpha\rVert_p^p\to0;
\]
$p=\infty$时把有限和换成有限个量的最大值，结论相同。
\end{proof}

证明的关键在第二步：测试恒等式两边都是对$L^p$收敛连续的线性量，所以约束在极限下保持。证明没有用到区域边界的形状，也没有得到逐点收敛或边界值，这些问题后面另行处理。

\subsection{\texorpdfstring{$H^k$}{Hk}空间}
\label{sec:sobolev-hk}

$p=2$时，函数与各阶导数都用平方积分度量，这些平方积分可以由内积给出。记$H^k(\Omega)=W^{k,2}(\Omega)$，定义
\begin{equation}
\langle f,g\rangle_{H^k}=\sum_{|\alpha|\leq k}\int_\Omega\partial^\alpha f\,\partial^\alpha g\,\mathrm dx .
\label{eq:sobolev-hk-inner}
\end{equation}

每一项都是$L^2(\Omega)$的内积，所以式~\eqref{eq:sobolev-hk-inner}有限，对两个变量都线性，并且对称。$f$与自身的内积正是Sobolev范数的平方，零阶项保证它只在$f=0$时为零。$k=1$时，
\[
\langle f,g\rangle_{H^1}=\int_\Omega fg\,\mathrm dx+\int_\Omega\nabla f\cdot\nabla g\,\mathrm dx .
\]

有限个$L^2$内积相加得到的只是内积空间。称它为Hilbert空间，依靠的是定理~\ref{thm:sobolev-complete}的完备性。有了Hilbert结构，第~\ref{chap:function-inner-product-spaces}章的正交投影就可以用于$H^k(\Omega)$的闭子空间，例如第~\ref{sec:sobolev-constrained-subspaces}节的零边界子空间。

\paragraph{一维$H^1$的加权坐标}
现在把“加权范数”的说法落实为定理。在区间$(0,1)$上，取式~\eqref{eq:sobolev-cosine-basis}的余弦基，并取另一组正交归一基$s_n(x)=\sqrt2\sin(\pi nx)$，$n\geq1$。两组基用求导互相联系：$e_n'=-\pi n\,s_n$。下面的定理说，$H^1(0,1)$恰好是坐标列的一个加权$\ell^2$空间。

\begin{theorem}[$H^1(0,1)$的余弦坐标]
\label{thm:sobolev-cosine-coordinates}
设$f\in L^2(0,1)$的余弦坐标为$a_n=\int_0^1fe_n\,\mathrm dx$。则$f\in H^1(0,1)$当且仅当$\sum_{n\geq0}\pi^2n^2a_n^2<\infty$，并且此时
\begin{equation}
\lVert f\rVert_{H^1}^2=\sum_{n\geq0}(1+\pi^2n^2)\,a_n^2 .
\label{eq:sobolev-cosine-norm}
\end{equation}
\end{theorem}

\begin{proof}
先设$f\in H^1(0,1)$，并取它的绝对连续代表（第~\ref{chap:smoothing-weak-derivatives}章）。乘积$fs_n$绝对连续，几乎处处有$(fs_n)'=f's_n+fs_n'$，而$s_n'=\pi n\,e_n$。在$[0,1]$上积分：
\[
\int_0^1f's_n\,\mathrm dx+\pi n\int_0^1fe_n\,\mathrm dx
=\bigl[fs_n\bigr]_0^1=0 .
\]
边界项为零，是因为$s_n(0)=s_n(1)=0$。所以$f'$的正弦坐标是$-\pi n\,a_n$。对$f'$用正弦基的Parseval等式，得到
\[
\lVert f'\rVert_2^2=\sum_{n\geq1}\pi^2n^2a_n^2 .
\]

反过来，设加权和有限。令部分和$F_N=\sum_{n\leq N}a_ne_n$，它是光滑函数，$F_N'=-\sum_{n\leq N}\pi n\,a_n\,s_n$。对$M<N$，
\[
\lVert F_N-F_M\rVert_{H^1}^2=\sum_{M<n\leq N}(1+\pi^2n^2)\,a_n^2\to0,
\]
所以$(F_N)$是$H^1(0,1)$中的Cauchy列。由定理~\ref{thm:sobolev-complete}它有极限，而$F_N\to f$在$L^2$中成立，所以极限就是$f$，$f\in H^1(0,1)$。式~\eqref{eq:sobolev-cosine-norm}由两次Parseval等式相加得到。
\end{proof}

余弦基的好处在证明里看得很清楚：分部积分的边界项被$s_n$在端点的零值消掉，所以定理对$H^1(0,1)$的每个函数都成立，不需要边界条件。可以用$f(x)=x$核对。分部积分得$a_n=\sqrt2((-1)^n-1)/(\pi^2n^2)$，只有奇数项非零，权重后的每项是$8/(\pi^2n^2)$，于是
\[
\sum_{n\text{奇}}\pi^2n^2a_n^2=\frac8{\pi^2}\cdot\frac{\pi^2}8=1=\lVert f'\rVert_2^2.
\]

把类比的对应关系和失效之处列在表~\ref{tab:sobolev-weight-analogy}中。最重要的是后两行：加权坐标只是$p=2$、$k=1$时最干净的图景，它帮助理解机制，但不能替代后面对一般$p$和一般区域的证明。

\begin{BookInlineTable}
  \centering
  \caption{有限维二次型与$H^1(0,1)$的对应。}
  \label{tab:sobolev-weight-analogy}
  \begin{tabular}{lll}
    \toprule
    & 三点差分 & $H^1(0,1)$\\
    \midrule
    坐标方向 & 常向量、斜坡、锯齿 & $e_0,e_1,e_2,\ldots$\\
    权 & $1,2,4$，有界 & $1+\pi^2n^2$，无界\\
    零导数方向 & 常向量 & 常数$e_0$\\
    $k\geq2$ & 权换成$D^\top D$的幂 & 需要边界相容条件\\
    $p\neq2$ & 没有正交坐标 & 没有正交坐标\\
    \bottomrule
  \end{tabular}
\end{BookInlineTable}

表中$k\geq2$一行需要解释。权$n^4$不再描述整个$H^2(0,1)$：$f(x)=x$属于$H^2(0,1)$，它的坐标却满足$n^4a_n^2\to$正常数（沿奇数$n$），加权和发散。分部积分此时要把一次导数的边界值也消掉，余弦坐标只描述满足$f'(0)=f'(1)=0$的子空间。周期函数没有这个问题，第~\ref{sec:sobolev-continuous-representatives}小节末尾还会用到周期情形。
